\input{preamble}

\begin{document}

\title{Ringed Spaces}
\maketitle

\phantomsection
\label{section-phantom}

\noindent
So far
this chapter is completely AI-generated.

\tableofcontents

\section{Motivation}
\label{section-motivation}

\medskip\noindent
Many geometric spaces
carry a sheaf of functions.
This records not only the
topology of the space,
but also the local algebra
of functions on it.

\medskip\noindent
The basic examples are
smooth manifolds with
smooth functions,
complex manifolds with
holomorphic functions,
affine varieties with
regular functions,
and schemes with their
structure sheaf.

\section{Presheaves}
\label{section-presheaves}

\begin{definition}
\label{definition-presheaf}
Let $X$ be a topological space.
A presheaf of sets $\mathcal F$
on $X$ assigns to each open
$U \subset X$ a set
$\mathcal F(U)$,
and to each inclusion
$V \subset U$
a restriction map
$$
\rho_{UV}:\mathcal F(U)\to\mathcal F(V).
$$
These maps satisfy
$$
\rho_{UU}=\mathrm{id}
$$
and for $W \subset V \subset U$,
$$
\rho_{UW}=\rho_{VW}\circ\rho_{UV}.
$$
\end{definition}

\begin{definition}
\label{definition-presheaf-abelian}
A presheaf of abelian groups,
rings, or modules is defined
in the same way,
with each $\mathcal F(U)$
an abelian group, ring,
or module, and each
restriction map respecting
the structure.
\end{definition}

\begin{remark}
\label{remark-restrictions}
The maps $\rho_{UV}$ are
called restriction maps.
If $s \in \mathcal F(U)$,
we often write $s|_V$
for $\rho_{UV}(s)$.
\end{remark}

\section{Sheaves}
\label{section-sheaves}

\begin{definition}
\label{definition-sheaf}
A presheaf $\mathcal F$
on $X$ is a sheaf if for every
open cover $U=\bigcup U_i$:

if $s,t \in \mathcal F(U)$
and $s|_{U_i}=t|_{U_i}$ for all
$i$, then $s=t$.

if $s_i \in \mathcal F(U_i)$
agree on overlaps
$U_i \cap U_j$,
then there exists a unique
$s \in \mathcal F(U)$
with $s|_{U_i}=s_i$.
\end{definition}

\begin{remark}
\label{remark-sheaf-intuition}
A sheaf is a presheaf whose
sections are determined locally
and can be glued from
compatible local data.
\end{remark}

\begin{example}
\label{example-sheaves}
The sheaves of continuous,
smooth, holomorphic, and
regular functions are
basic examples.
\end{example}

\section{Stalks and germs}
\label{section-stalks}

\begin{definition}
\label{definition-stalk}
Let $\mathcal F$ be a sheaf on
$X$ and let $p \in X$.
The stalk $\mathcal F_p$
is the set of germs of sections
near $p$.
An element is represented by
a pair $(U,s)$ with
$p \in U$ and $s \in \mathcal F(U)$.
We declare
$(U,s)\sim(V,t)$ if there is a
neighborhood $W \subset U \cap V$
of $p$ such that
$s|_W=t|_W$.
\end{definition}

\begin{theorem}
\label{theorem-stalk-iso}
Let $\varphi:\mathcal F\to\mathcal G$
be a morphism of sheaves.
Then $\varphi$ is an isomorphism
if and only if
each induced map
$\mathcal F_p\to\mathcal G_p$
is an isomorphism.
\end{theorem}

\begin{proof}
If $\varphi$ is an isomorphism,
then each stalk map is an
isomorphism by functoriality.

Conversely, assume all stalk maps
are isomorphisms.
To prove injectivity, let
$s \in \mathcal F(U)$ satisfy
$\varphi(s)=0$.
For each $p \in U$,
the germ of $s$ at $p$
maps to $0$ in $\mathcal G_p$,
so the germ of $s$ is $0$
in $\mathcal F_p$.
Hence there is a neighborhood
where $s$ vanishes.
By the local identity axiom,
$s=0$.

For surjectivity, let
$t \in \mathcal G(U)$.
For each $p \in U$,
choose a germ $s_p \in \mathcal F_p$
mapping to the germ of $t$.
Represent $s_p$ on some
neighborhood $U_p$.
The family of local lifts
agrees on overlaps because the
stalk maps are injective.
By gluing, these local lifts
come from a section of
$\mathcal F(U)$ mapping to $t$.
\end{proof}

\section{Kernels and cokernels}
\label{section-kernels-cokernels}

\begin{definition}
\label{definition-kernel-image-cokernel}
Let $\varphi:\mathcal F\to\mathcal G$
be a morphism of sheaves of
abelian groups or modules.
The kernel is defined on opens by
$$
(\ker\varphi)(U)=\ker(\mathcal F(U)\to
\mathcal G(U)).
$$
The image and cokernel are
defined as presheaves and then
sheafified when needed.
\end{definition}

\begin{remark}
\label{remark-image-cokernel}
Kernels are computed directly
on open sets.
Images and cokernels usually
must be sheafified.
A quotient of sheaves of rings
is usually not a sheaf of rings;
it is naturally a sheaf of
abelian groups or modules.
\end{remark}

\section{Exact sequences}
\label{section-exact-sequences}

\begin{definition}
\label{definition-exact-stalkwise}
A sequence
$$
0\to \mathcal F'\to \mathcal F
\to \mathcal F''\to 0
$$
of sheaves of abelian groups or
modules is exact if the induced
sequence on every stalk is exact.
\end{definition}

\begin{proposition}
\label{proposition-exact-stalks}
Exactness of sheaves can be checked
on stalks.
\end{proposition}

\begin{proof}
For sheaves of abelian groups or
modules, exactness on opens
implies exactness on stalks
by passing to germs.
Conversely, if a morphism has
zero stalk at every point,
then the section vanishes on a
cover by neighborhoods, hence
vanishes globally by the
sheaf axiom.
This gives exactness at each
term.
\end{proof}

\begin{remark}
\label{remark-exact-sequences-use}
Exact sequences of sheaves are
the basic tool for comparing
local and global invariants.
\end{remark}

\section{Pushforward sheaves}
\label{section-pushforward}

\begin{definition}
\label{definition-pushforward}
Let $f:X\to Y$ be continuous and
let $\mathcal F$ be a sheaf on
$X$.
The pushforward sheaf is
defined by
$$
(f_*\mathcal F)(U)
=
\mathcal F(f^{-1}(U)).
$$
\end{definition}

\begin{remark}
\label{remark-pushforward-sheaf}
This is a sheaf on $Y$.
It lets us view a sheaf upstairs
downstairs by assigning to each
open set in $Y$ the sections
over its inverse image.
\end{remark}

\section{Pullback sheaves}
\label{section-pullback}

\begin{definition}
\label{definition-pullback}
Let $f:X\to Y$ be continuous and
let $\mathcal G$ be a sheaf on
$Y$.
Informally, the inverse image
sheaf $f^{-1}\mathcal G$
is the sheafification of the
presheaf
$$
U\mapsto
\varinjlim_{f(U)\subset V}
\mathcal G(V).
$$
For sheaves of modules one
usually sets
$$
f^*\mathcal G
=
f^{-1}\mathcal G
\otimes_{f^{-1}\mathcal O_Y}
\mathcal O_X.
$$
\end{definition}

\section{Rings and modules}
\label{section-rings-modules}

\begin{definition}
\label{definition-sheaf-ring}
A sheaf of rings is a sheaf
whose sections over each open
set form a ring.
An $\mathcal O_X$-module is a
sheaf of abelian groups
with an action of the sheaf of
rings $\mathcal O_X$.
\end{definition}

\begin{remark}
\label{remark-module-example}
If $(X,\mathcal O_X)$ is a
ringed space, then an
$\mathcal O_X$-module is a sheaf
locally acted on by functions.
\end{remark}

\section{Ringed spaces}
\label{section-ringed-spaces}

\begin{definition}
\label{definition-ringed-space}
A ringed space is a pair
$$
(X,\mathcal O_X)
$$
where $X$ is a topological space
and $\mathcal O_X$ is a sheaf
of rings on $X$.
\end{definition}

\begin{example}
\label{example-ringed-spaces}
Smooth manifolds, complex
manifolds, affine varieties,
and schemes are ringed spaces.
\end{example}

\section{Locally ringed spaces}
\label{section-locally-ringed}

\begin{definition}
\label{definition-locally-ringed}
A ringed space $(X,\mathcal O_X)$
is locally ringed if each stalk
$\mathcal O_{X,p}$ is a local
ring.
\end{definition}

\begin{remark}
\label{remark-local-natural}
This is geometrically natural:
the maximal ideal in
$\mathcal O_{X,p}$ consists of
functions vanishing at $p$.
\end{remark}

\begin{definition}
\label{definition-lrs-morphism}
A morphism of locally ringed
spaces is a morphism of ringed
spaces whose induced maps on
stalks are local homomorphisms.
\end{definition}

\section{Gluing sheaves}
\label{section-gluing-sheaves}

\begin{lemma}
\label{lemma-gluing-sheaves}
Let $X=\bigcup U_i$ be an open
cover.
Suppose $\mathcal F_i$ are
sheaves on $U_i$ and the
restrictions to $U_i\cap U_j$
are identified by isomorphisms
satisfying the cocycle
condition.
Then there is a unique sheaf
$\mathcal F$ on $X$ whose
restriction to $U_i$ is
$\mathcal F_i$.
\end{lemma}

\begin{proof}
Define $\mathcal F(U)$ to be
families of sections on the
pieces $U\cap U_i$
which agree on overlaps.
The sheaf axioms follow from
those of the $\mathcal F_i$.
Uniqueness is built into the
construction.
\end{proof}

\section{Gluing ringed spaces}
\label{section-gluing-ringed-spaces}

\medskip\noindent
Ringed spaces can be glued by
gluing both the underlying spaces
and the structure sheaves.
This is one of the basic ways to
build global objects from local
charts.

\section{Affine schemes}
\label{section-affine-schemes}

\medskip\noindent
Very briefly, $\Spec A$ carries a
structure sheaf.
For a basic open set $D(f)$,
$$
\mathcal O_{\Spec A}(D(f))=A_f.
$$

\section{Complex spaces}
\label{section-complex-spaces}

\medskip\noindent
Complex spaces are locally ringed
spaces locally modeled on
quotients of sheaves of
holomorphic functions.
We do not develop complex
analytic geometry here.

\section{Skyscraper sheaves}
\label{section-skyscraper-sheaves}

\begin{definition}
\label{definition-skyscraper}
Let $p\in X$ and let $A$ be an
abelian group.
The skyscraper sheaf at $p$
with value $A$ is the sheaf
whose sections on an open set
$U$ are $A$ if $p\in U$ and $0$
otherwise.
For finitely many points one
takes the direct sum of the
corresponding skyscraper sheaves.
\end{definition}

\begin{theorem}
\label{theorem-skyscraper-cohomology}
In the elementary cases needed
later, a skyscraper sheaf
supported on finitely many
points has no higher
cohomology.
\end{theorem}

\begin{proof}
Such a sheaf is a finite direct
sum of sheaves supported on
points, and each of these is
acyclic on the usual spaces
that will occur later, because
its sections are already
concentrated on a finite closed
set.
\end{proof}

\section{Cohomology and Euler characteristic}
\label{section-cohomology-euler}

\medskip\noindent
At a basic level, the
Euler characteristic of a sheaf
$\mathcal F$ is
$$
\chi(\mathcal F)
=
\sum_i (-1)^i h^i(X,\mathcal F).
$$

\begin{proposition}
\label{proposition-euler-additive}
Short exact sequences of sheaves
give long exact sequences in
cohomology.
Hence the Euler characteristic
is additive on short exact
sequences.
\end{proposition}

\begin{proof}
The long exact sequence in
cohomology is the standard
derived-functor statement.
Taking alternating sums of
dimensions in a long exact
sequence gives additivity of
Euler characteristic.
\end{proof}

\section{Finite maps}
\label{section-finite-maps}

\begin{theorem}
\label{theorem-finite-pushforward}
If $f:X\to Y$ is finite and
$\mathcal F$ is coherent on $X$,
then
$$
H^i(Y,f_*\mathcal F)
\cong
H^i(X,\mathcal F).
$$
\end{theorem}

\begin{proof}
Finite maps are affine.
Affine maps have no higher
direct images for
quasi-coherent sheaves.
The Leray spectral sequence
then gives the claimed
identification.
\end{proof}

\input{chapters}
\bibliography{my}
\bibliographystyle{amsalpha}
\end{document}
